% ============================================================================
% 4. BASELINES
% ============================================================================
\section{Baselines}

We implement four baselines spanning the text-to-graph-augmented spectrum,
plus an agentic Tool-Use baseline:

\paragraph{Flat Text (FT).}
The complete schema is serialized as a text document listing all tables,
FK relationships, and lineage edges. This forms the strongest
``memorize everything'' baseline, as the LLM receives the complete graph.

\paragraph{Vector RAG (VR).}
Tables are chunked with their local FK and lineage edges, embedded using
Sentence-BERT~\cite{reimers2019sentencebert}, and retrieved via cosine
similarity ($k=15$). A global lineage summary is appended to provide
cross-chunk edge visibility.

\paragraph{Graph-Augmented (GA).}
Starting from the tables referenced in the question, a 3-hop BFS
over the schema graph extracts the local subgraph neighborhood. The
extracted context includes: (i)~the list of tables within the BFS
frontier, (ii)~FK edges between extracted tables, (iii)~lineage
edges (\texttt{derived\_from}) between extracted tables, and
(iv)~per-table column lists. Crucially, GA does \emph{not} inject
algorithm outputs (shortest paths, component IDs, or hop distances)
into the context---the LLM must infer these from the raw subgraph.
This design ensures that GA tests \emph{structural reasoning} over
graph context rather than \emph{answer copy-paste} from precomputed
results. A learned graph encoder (e.g., GNN) represents future work.

\paragraph{Tool-Use (TU).}
An agentic baseline providing the LLM with 9 graph algorithm tools:
\texttt{shortest\_path}, \texttt{connected\_components},
\texttt{check\_fk\_adjacency}, \texttt{get\_lineage\_forward/reverse},
\texttt{transitive\_lineage}, \texttt{get\_component\_of},
\texttt{get\_fk\_neighbors}, and \texttt{list\_tables}. The LLM
receives a schema summary (table list + edge types) and tool
definitions, then engages in a multi-turn conversation (up to 3 tool
calls per question) implemented via function calling. This baseline
is closest in spirit to graph-grounded agent systems such as
ToG~\cite{sun2024tog} and GraphCoT~\cite{jin2024graphcot}, adapted
for schema topology reasoning.

All baselines are evaluated with two frontier models: Gemini~2.5 Flash
(closed-weight) and DeepSeek-V3 (open-weight, 671B MoE). We use greedy
decoding ($\tau{=}0.0$) with a single deterministic pass to eliminate
sampling variance. Responses use a structured JSON format with reasoning
and answer fields. Note that while our PyG \texttt{HeteroData} graphs
store 6 node-level structural features (in-degree, out-degree, PageRank,
betweenness centrality, etc.), \emph{none of these features are used by
any current baseline at inference time}. They are provided for future work
on learned graph encoders (e.g., GNNs).

\paragraph{Oracle.}
Injects gold graph algorithm outputs (the correct shortest path, connected
component membership, lineage traversal result, etc.) directly into the
prompt. Oracle scores represent the upper bound achievable by each
model---any remaining errors reflect decoding or format issues, not
retrieval or reasoning failures.

\paragraph{List canonicalization and tie handling.}
Unordered list answers are lowercased and sorted before comparison;
exact match requires identical sorted sets. For shortest-path questions,
multiple valid routes may exist (e.g., two equally short FK paths).
Our scoring function validates the \emph{predicted path} against the
FK adjacency matrix: if the predicted sequence of tables corresponds
to a valid path of the same length as the gold answer, it is marked
correct even if it differs from the gold path. This ensures EM is
robust to tie-breaking variability. Boolean and integer answers are
type-cast before comparison.
